% compile with lualatex, not pdflatex etc.

\documentclass{article}

\usepackage{amsmath, amssymb, amsthm}
\usepackage{luacode}

\directlua{require('piformulas.lua')}


\title{Evaluating Pi Formulas with Lua}

\author{Chris Webb, LaTeX Snippets}
\date{}


\begin{document}

\maketitle


\section{Hard Coded Return Value}


\[
\pi \approx 
\luaexec{tex.print(pi())}
\]


\section{Fraction}

\[
\frac{22}{7} \approx 
\luaexec{tex.print(fraction())}
\]


\section{Francois Viete}

\[
\frac{2}{\pi} = 
\frac{\sqrt{2}}{2} \cdot 
\frac{\sqrt{2+\sqrt{2}}}{2} \cdot
\frac{\sqrt{2+\sqrt{2+\sqrt{2}}}}{2} 
\cdots
\]

\[
\pi \approx 
\luaexec{tex.print(viete(16))}
\;\text{(16 terms)}
\]


\section{John Wallis}

\[
\frac{\pi}{2} = 
\frac{2}{1} \cdot
\frac{2}{3} \cdot
\frac{4}{3} \cdot
\frac{4}{5} \cdot
\frac{6}{5} \cdot
\frac{6}{7} \cdot
\frac{8}{7} \cdot
\frac{8}{9}
\cdots
\]

\[
\pi \approx 
\luaexec{tex.print(wallis(100000))}
\;\text{(100000 terms)}
\]


\section{John Machin}

\[
\frac{\pi}{4} = 4 \arctan \frac{1}{5} - \arctan \frac{1}{239}
\]

\[
\pi \approx 
\luaexec{tex.print(machin())}
\]


\section{Gregory-Leibniz}

\[
\pi = 
\frac{4}{1} -
\frac{4}{3} +
\frac{4}{5} -
\frac{4}{7} +
\frac{4}{9} -
\frac{4}{11} +
\frac{4}{13} -
\cdots
\]

\[
\pi \approx 
\luaexec{tex.print(gregoryleibniz(100000))}
\;\text{(100000 terms)}
\]


\section{Nilakantha}

\[
\pi = 3 +
\frac{4}{2 \times 3 \times 4} -
\frac{4}{4 \times 5 \times 6} +
\frac{4}{6 \times 7 \times 8} -
\frac{4}{8 \times 9 \times 10} +
\cdots
\]

\[
\pi \approx 
\luaexec{tex.print(nilakantha(100000))}
\;\text{(100000 terms)}
\]


\end{document}